% ARCHIVED DRAFT: independently audited by root and incorporated on 2026-11-19 UTC.
% The assembled paper is authoritative; this file retains drafting provenance.
\begin{saresult}{B.075}{Proved}{Infinite-residue conductor cusps with arbitrary coefficient extensions}
Let $k=\mathbb F_q\subset l$ with $l$ an infinite field, let $a\ge1$, and put
\[
 V=l[[t]],\qquad K=l((t)),\qquad R=k+t^aV,\qquad
                         \mathfrak m=t^aV.
\]
For every unital commutative inclusion $R\subset S$, with arbitrary
$R$-torsion in $S$, the following are equivalent:
\[
\begin{gathered}
 \text{the polynomial co-operations exist uniquely};\\
 \theta_2\text{ is injective};\qquad
 (R:_S\mathfrak m)=R;\\
 \operatorname{Ann}_{S/R}(\mathfrak m)=0;\qquad
 \Omega(R,S)=R[X].                                       \tag{B75.1}
\end{gathered}
\]
Every positive case has all multivariable tensor isomorphisms.
This removes the torsion-free restriction from the entire family B.043,
including the explicit base $\mathbb F_2+t\mathbb F_2(u)[[t]]$.
\end{saresult}
\begin{proof}
We first identify the coefficient ring without any localization theorem
for integer-valued modules. The determinant ideals $D_d(R)$ are cofinal
with the powers of $\mathfrak m$. For fixed $d$, choose $d+1$ elements
of $t^a l$ with distinct coefficients in $l$. Their determinant has
the form $t^{ad(d+1)/2}u$ with $u\in V^\times$. Multiplying it by
$\mathfrak m$ shows that $D_d(R)$ contains a power of $\mathfrak m$.
In the opposite direction, every tuple of $d+1$ nodes has residue
multiplicities $n_1,\ldots,n_q$; its determinant belongs to
$\mathfrak m^{\sum_i\binom{n_i}{2}}$. The minimum of this exponent
tends to infinity with $d$. Thus B.064 gives the actual coefficient ring
\[
 C=\bigcup_{n\ge0}(R:_S\mathfrak m^n),\qquad
                         \Omega(R,S)=\Omega(R,C).         \tag{B75.2}
\]
It is a ring. If $(R:_S\mathfrak m)=R$, iterating colons gives $C=R$,
so every positive conclusion follows. We prove that tensor injectivity
forces $C=R$.

Assume $\theta_2$ is injective, so $R\to C$ is an epimorphism by B.064.
Localizing at the nonzero elements of $R$ maps every element of $C$ into
the embedded fraction field $K$: some nonzero element of
$\mathfrak m^n$ multiplies it into $R$. Let $L\subset K$ be this image,
and $T$ its kernel. Every element of $T$ is annihilated by a power of
$\mathfrak m$, because its products with that power lie in $R$ and
vanish after localization. The composite $R\to L$ is an epimorphism.

If $L\subset V$, the ideal $\mathfrak m$ is an ideal of $L$ with
$\mathfrak mL=\mathfrak m$. Reduction gives an injective epimorphism
from $k$ to $L/\mathfrak m$, forcing equality and $L=R$.
If $L$ contains an element of negative valuation, a positive power of
its inverse belongs to $\mathfrak m$. Thus $L$ contains the inverse of
some nonzero $b\in\mathfrak m$, and $bV\subset\mathfrak m\subset R$
implies $V\subset L$. An overring of the discrete valuation ring $V$
with an element of negative valuation is $K$. Hence $L=R$ or $K$.
In the former case rationalization is a ring retraction $C\to R$;
epicness forces $C=R$.

In the latter case $C$ has a precise elementary decomposition. Put
$s=t^a\in R$, choose $c\in C$ with rational image $s^{-1}$, and set
$x=sc$. Since $x-1\in T$, choose $N\ge1$ with $s^N(x-1)=0$.
Multiplication by $c^N$ gives $x^N(x-1)=0$. Therefore
$e=x^N$ is an idempotent and $ex=e$. On $eC$, the element $s$ is
invertible with inverse $ec$. All $s$-power torsion vanishes there,
so rationalization identifies $eC$ with $K$.
The unit $1-e$ of $D=(1-e)C$ belongs to $T$, so
$\mathfrak m^M D=0$ for some $M$.
The epimorphism $R\to D$ is surjective. Indeed, modulo its nilpotent
maximal-ideal image it is an epimorphism from the field $k$, hence an
isomorphism if $D\ne0$. Thus $D$ is the image of $R$ plus
$\mathfrak mD$, and finite iteration of $\mathfrak m^MD=0$ proves
surjectivity. Consequently
\[
             C\simeq K\times R/I,\qquad
             \mathfrak m^M\subset I,                     \tag{B75.3}
\]
with the diagonal inclusion of $R$; $I=R$ is allowed and means $C=K$.

We exclude all possibilities in (B75.3) by nonzero tensor kernels.
Write $P(X)=X^q-X$ and $U=P(X)/t^a$. The following explicit ring is
central to the argument:
\[
 B=\{F\in K[X]:F(R)\subset V\}
   =V[X,U]/(X^q-X-t^aU)
   =\bigoplus_{0\le i<q,\ j\ge0}VX^iU^j.                 \tag{B75.4}
\]
Here the presentation means its embedded image in $K[X]$.
The displayed monomials are a $K$-basis of $K[X]$, by successive
division by the monic $P$. They all preserve $R$ into $V$.
If an integer-valued polynomial had a coefficient of negative minimum
$t$-valuation in this basis, multiplication by the opposite power of
$t$ would give a polynomial with nonzero reduction modulo $t$ whose
values are all zero modulo $t$. But
\[
 B/tB\simeq l[X,U]/(X^q-X)\simeq\prod_{b\in k}l[U],       \tag{B75.5}
\]
and evaluating $X=b+t^av$ with $v\in l$ sends $U$ modulo $t$ to $-v$.
The infinitude of $l$ makes such a nonzero polynomial impossible.
This proves (B75.4).

For $I=R$, the same evaluation modulo $t^a$ gives the split module
formula
\[
 \Omega(R,K)=
 \bigoplus_{0\le i<q}RX^i
 \ \oplus\!!\bigoplus_{0\le i<q,\ j\ge1}t^aVX^iU^j.     \tag{B75.6}
\]
To justify all $a$ layers, on constant $v\in l$ one has
$X\equiv b$ and $U\equiv-v$ modulo $t^a$.
The coefficient of $t^0$ must be a $k$-valued polynomial on the
infinite field $l$, hence constant in $k$; every coefficient of
$t^r$, $0<r<a$, must vanish. Interpolation on the $q$ possible $b$
then gives exactly (B75.6).
Choose $\alpha\in l\setminus k$. The tensor
\[
                  \alpha P\otimes P-P\otimes\alpha P     \tag{B75.7}
\]
has zero polynomial image and is nonzero. Projection to the
$(i,j)=(0,1)$ summand of (B75.6), followed by division by $t^a$ and
reduction modulo $t$, sends it to
$\alpha\otimes1-1\otimes\alpha\ne0$ in $l\otimes_k l$.
The projection is $R$-linear via $R\to k$. This excludes $C=K$.

Now suppose $I\ne R$. Choose the least valuation $b\ge a$ of a
nonzero element of $I$. Multiplying such an element by
$\mathfrak m=t^aV$ gives
\[
                      t^{b+a}V\subset I\subset t^bV.
\]
Put $H=\{F\in K[X]:F(R)\subset I\}$ and let $E$ be the first projection
of $\Omega(R,K\times R/I)$ into $K[X]$. Lifting the coefficients of
its second polynomial component to $R$ shows exactly that
\[
                  E=R[X]+H,\qquad
                  t^{b+a}B\subset H\subset t^bB.           \tag{B75.8}
\]
We claim $H\not\subset\mathfrak mE$.

For a polynomial in $t^jB$, call its image after division by $t^j$
in (B75.5) its leading residue at level $j$.
Let $d$ be the least level for which some element of $H\cap t^dB$
has a leading residue with a nonconstant $U$-component. Such a level
exists with $b\le d\le b+a$, since $t^{b+a}U\in H$.
Choose an element witnessing it and a component indexed by
$b_0\in k$ whose residue polynomial $h(U)$ is nonconstant.
Precompose it with
\[
                    \psi_N(X)=b_0+t^aU(X)^N.
\]
This polynomial maps $R$ into $R$, so the composition remains in $H$.
Its leading residue at level $d$ is $h(-U^N)$ in every component,
since $\psi_N$ has residue $b_0$ and
$U(\psi_N)\equiv-U^N\pmod t$.
Taking $N$ large gives an element $F\in H\cap t^dB$ whose leading
$U$-degree is strictly larger than $\lfloor d/a\rfloor-1$.

This $F$ cannot lie in $\mathfrak mE$.
In the monic basis $X^iP^j$, one has
\[
 \mathfrak mR[X]=
       \bigoplus_{0\le i<q,\ j\ge0}t^{a(j+1)}VX^iU^j.
\]
Therefore the leading residue at level $d$ of an element of
$\mathfrak mR[X]\cap t^dB$ has $U$-degree at most
$\lfloor d/a\rfloor-1$.
Also $\mathfrak mH\subset t^{b+a}B$. If $d<b+a$ it contributes
nothing at level $d$. If $d=b+a$, its leading residues are linear
combinations over $l$ of leading residues of $H$ at level $b$;
these are $U$-constant by the minimality of $d$.
Thus adding $\mathfrak mH$ cannot supply the missing higher degree.
This proves $F\notin\mathfrak mE$.

The polynomial $f=(F,0)$ belongs to $\Omega(R,C)$, and is nonzero
modulo $\mathfrak m\Omega$ because its first projection is outside
$\mathfrak mE$. Choose $n$ with $an\ge b+a$ and put
$g=(0,P^n\bmod I)$. It is a null polynomial on $R$, so also belongs
to $\Omega(R,C)$. Its second component is monic and remains nonzero
over $k=R/\mathfrak m$, so $g\notin\mathfrak m\Omega$.
The pure tensor
\[
                              f\otimes g                  \tag{B75.9}
\]
has zero polynomial image, because its two coefficient components
have disjoint idempotent support. It is nonzero: reduction modulo
$\mathfrak m$ gives the tensor of two nonzero vectors over the field
$k$. This excludes every proper $I$ in (B75.3).

All cases with $C\ne R$ have contradicted tensor injectivity.
Finally $C=R$ is equivalent to $(R:_S\mathfrak m)=R$ by the iterated
colon identity, and the latter is precisely the annihilator condition
in (B75.1). B.064 supplies the remaining conclusions.
\end{proof}
